% Canonical Paper IV section.
Arithmetic evaluation is ordinarily described by a map from an expression to
its value.  Arithmetic Expression Geometry begins one level earlier.  It treats
an ordered evaluation history as mathematical data and asks what survives when
that history is replaced successively by an operator, a quotient state, or an
endpoint.  Paper~I supplies the marked-history formalism, projective evaluation,
the affine sector, and the ACS charge shadow \cite{Yuan2026AEGFoundations}.
Paper~II develops analysis on the regular affine spaces
\cite{Yuan2026AEGFunctionTheory}, and Paper~III studies singular zero geometry
and topological transport \cite{Yuan2026AEGZeroGeometry}.  The present paper
returns to the forgetful maps themselves.

The motivating hierarchy is
\begin{equation}
  \text{marked history}
  \longrightarrow
  \text{projective operator}
  \longrightarrow
  \text{quotient state}
  \longrightarrow
  \text{endpoint or observation}.
  \label{eq:intro-hierarchy}
\end{equation}
Every arrow can identify distinct inputs.  It is tempting to call the resulting
fiber ``lost information'' and to identify its logarithmic size with memory,
representation complexity, or reconstruction time.  Such an identification is
not valid without further data.  A fiber is an algebraic or set-theoretic
object.  Computational complexity belongs to an encoded family of tasks in a
specified machine model.  The missing bridge is a theorem saying which
distinctions in a fiber can still affect an allowed future.

\subsection{Four layers}

We organize the subject into four layers.
\begin{enumerate}
  \item \textbf{Projective quotient fibers.}
  Bilateral arithmetic histories evaluate in $G=\PGL_2(K)$.  Homogeneous
  quotients such as $G/H$ forget a projective frame direction while retaining
  an image--kernel projector.

  \item \textbf{Contextual residuals.}
  At a typed cut $C$, two pasts are equivalent when every common legal future
  gives the same declared observation.  This equivalence depends on the future
  class, the observation, and the treatment of undefined arithmetic.

  \item \textbf{Operational live configurations.}
  A run stores concrete values, checkpoints, messages, or clauses.  Its space
  is the encoded size of what is live, not the size of an abstract quotient
  fiber and not the set of nodes that have ever been computed.

  \item \textbf{Rewrite fibers.}
  Equal semantic objects may have many factorizations, parenthesizations, or
  schedules.  Rewrite distance, normal-form cost, filling area, and possible
  holonomy are properties of this realization fiber, not coordinates of an
  endpoint alone.
\end{enumerate}

Figure~\ref{fig:intro-four-layers} records the direction of comparison used in
this paper.  The vertical arrows are not asserted to be isomorphisms.

\begin{figure}[ht]
\centering
\begin{tikzpicture}[
  node distance=7mm and 13mm,
  box/.style={draw,rounded corners,align=center,minimum width=38mm,minimum height=10mm},
  arr/.style={-{Latex[length=2.2mm]},thick}
]
\node[box] (hist) {marked histories\\and operators};
\node[box,right=of hist] (quot) {projective\\quotient fibers};
\node[box,below=of quot] (res) {contextual\\residual classes};
\node[box,left=of res] (live) {operational\\live configurations};
\node[box,below=of live] (rew) {rewrite fibers\\and schedules};
\draw[arr] (hist) -- node[above]{evaluation / quotient} (quot);
\draw[arr] (quot) -- node[right,align=left]{sufficiency requires\\a compatible future action} (res);
\draw[arr] (res) -- node[above,sloped]{encoding and realization} (live);
\draw[arr] (rew) -- node[left]{choose a run} (live);
\draw[arr,dashed] (hist) -- (rew);
\end{tikzpicture}
\caption{The four layers of Paper~IV.  A projective quotient becomes a
future-sensitive state only after the future action and observation are fixed;
a residual lower bound becomes operational space only after encoding and
machine conventions are fixed.}
\label{fig:intro-four-layers}
\end{figure}

\subsection{Main results}

The paper establishes six groups of results.

First, we complete the algebraic quotient layer.  A regular bivaluation---an
ordered pair of distinct projective points---is naturally a rank-one
idempotent.  For the ordered-pair stabilizer $H$ and the two point stabilizers
$B_\pm$,
\[
  G/H\cong\Biv_K^{\mathrm{reg}},
  \qquad
  G/B_\pm\cong\PP^1(K).
\]
The principal $H$-torsor $G\to G/H$ is the projective-frame residue over a
condensed projector.  This is algebraic condensation.  It must be distinguished
from convergence of normalized matrix products to a singular rank-one map.

Second, we prove a continuation-side theorem.  Left multiplication by every
future operator descends to $G/H$.  A single right update by $k$ descends
exactly when $k^{-1}Hk\subseteq H$; a reversible right action descends exactly
when $k$ normalizes $H$.  Consequently, the same quotient may be sufficient
for chronological continuation by left composition and insufficient for a
right-composition convention.  The theorem gives an exact meaning to the
statement that forgotten frame data can become visible again.

Third, we define contextual continuation residuals and prove the finite exact
minimal-state theorem.  The quotient by residual equivalence is a canonical
deterministic online state space, in the same structural tradition as the
Nerode equivalence for automata \cite{Nerode1958}.  Any exact implementation
must inject this quotient into its encoded states.  This yields a bit lower
bound and, under an explicit fixed-register snapshot convention, a
capacity--time lower bound.

Fourth, we introduce a live-configuration cost vector separating static
description, work, elapsed steps, causal depth, peak workspace, memory--time,
and communication.  We prove that a monotone completed-node set cannot encode
space in models allowing deletion or recomputation.

Fifth, we give a word-fiber and conditional-entropy inequality.  Exponential
raw history growth combined with small evaluation-image growth forces large
operator fibers and large conditional information.  The conclusion is about
an evaluation map.  It is not a time lower bound.  Finite generating sets give
quasi-isometric word metrics, so polynomial versus exponential group growth is
stable under that restricted change of generators, while computational costs
remain model-dependent.

Sixth, exact case studies calibrate the distinctions.  Horner words exhibit
charge condensation; OBDDs exhibit future-residual sensitivity to variable
order; butterflies and NTTs exhibit representation change and static torus
labels; matrix chains and checkpointing exhibit equal semantics with different
work and space; and rewrite potentials exhibit a telescoping obstruction to
nontrivial holonomy.

\subsection{What is not claimed}

No result in this paper asserts any of the implications
\[
  \text{noncommutativity}\Longrightarrow\text{negative curvature},
  \qquad
  \text{negative curvature}\Longrightarrow\text{exponential time},
\]
\[
  \text{exponential word growth}\Longrightarrow\text{NP-hardness},
  \qquad
  \log|\text{fiber}|=\text{workspace}.
\]
The free group already has exponential metric growth while free reduction is
linear in the input-word length.  Conversely, the OBDD example in
Section~\ref{sec:obdd} has an exponential cut-width separation although the
underlying variable restrictions commute.  The appropriate conclusions are
therefore conditional comparison theorems, not a universal scalar complexity
law.

\subsection{Roadmap}

Sections~\ref{sec:history-equality}--\ref{sec:projective-quotients} establish
the history and homogeneous-space layers.  Section~\ref{sec:condensation-continuation}
proves the quotient-continuation results.  Section~\ref{sec:contextual-residuals}
develops residual semantics, and Section~\ref{sec:operational-resources}
separates it from operational resources.  Sections~\ref{sec:growth-information}
and \ref{sec:rewrite-fibers} treat growth and rewriting.  The four case-study
sections follow.  The final two sections state the claim boundary and the
remaining research programme.  Appendices collect calculations, proof variants,
finite counts, and a theorem-level claim ledger.
